
\section{Compression et Distribution}

\subsection{Compression de fichiers inverse}

Soit un fichier inverse contenant 100 000 termes diff�rents, et 100 millions de documents. Chaque terme contient en moyenne 500 000 documents (peut varier de 10 000 � 5 millions). Nous nous int�ressons maintenant � la taille prise par cet index, et la mani�re de le compresser.

Pour cela, il faut savoir que :
\begin{itemize}
\item Un terme est cod� maximum sur 20 octets (avec une moyenne de 8 octets)
\item Un pointeur sur 8 octets
\item Un nombre est cod� sur 4 octets
\end{itemize}
Un index contient un $terme$, la \textit{fr�quence} du terme et un pointeur sur la \textit{posting list}.

\subsubsection{Compression de l'index de termes}

\begin{enumerate}
\item Quelle est la taille de l'index des termes ? (sans les \textit{posting lists})
	\begin{correction}
		100 000 x (20 + 4 + 8) = 3 200 000 o = 3,20 Mo
	\end{correction}

\begin{comment}
	\item Proposer une technique de compression des termes. Donner la taille de cet index.
	%\begin{comment}
	\begin{correction}
		En codant tous les termes sur leur taille moyenne, puis en rajoutant un pointeur sur la chaine enti�re � part (si besoin), nous pouvons r�duire la taille � 8o + 8o pour chaque terme. Ce qui nous fait descendre � 2,4 Mo.
	\end{correction}
	%\end{comment}

	\item Remplacer le codage du terme par un $termId$ (identifiant de terme dans le dictionnaire). Donner la taille de cet index.
	%\begin{comment}
	\begin{correction}
		Chaque triplet est alors cod� sur : (4 + 4 + 8).\\
		Ce qui nous fait une taille totale de 1,6 Mo.
	\end{correction}
	%\end{comment}

	\item Que pourra-t'on utiliser pour retrouver le terme en fonction de son $termId$ ? ?
	%\begin{comment}
	\begin{correction}
		On pourra utiliser un BTree pour retrouver chaque terme.
	\end{correction}
	%\end{comment}
\end{comment}
\end{enumerate}

\subsubsection{Compression de la \textit{Posting List}}

\begin{enumerate}
	\item Compte tenu du nombre de documents, donner la taille d'encodage d'un $docId$.
	\begin{correction}
		log$_2$(10$^8$) ~ 26,57 =$>$ 4 octets
	\end{correction}

	\item Donner, � partir de la taille moyenne des \textit{posting lists} et de la taille d'un $docId$, la taille d'un vecteur de l'index. Puis la taille de l'index global.
	\begin{correction}
		500 000 x 4 = 2 Mo / vecteur.\\
		100 000 x 2M = 200 Go.
	\end{correction}

	\item Proposer une solution pour encoder les docId d'une mani�re diff�rente. Donner le principe de l'encodage.
	\begin{correction}
		Une solution est de calculer les distances entre chaque docId avec des encodages � taille variable. Pour cela, � chaque docId, on calcul la diff�rence avec le pr�c�dent et on l'encode sur le nombre de bits n�cessaire. Le premier bit de chaque octet permet de dire si c'est la fin de l'encodage, s'il est � 1 c'est que c'est la fin du calcul, sinon on ajoute avec la suite.\\

	\end{correction}

	\item  Donner la transformation, puis l'encodage de la \textit{posting list} suivante : 105, 117, 222, 702, 3002
	\begin{correction}
	Transformation : 105, 2, 115, 480, 2300\\
	Pour chaque diff�rence : 1101001, 10, 1110011, 11 1100000, 10001 1111100\\
	Encodage : \textbf{1}1101001\textbf{1}0000010\textbf{1}1110011\textbf{0}0000011\textbf{1}1100000\textbf{0}0010001\textbf{1}1111100
	\end{correction}

	\item Donner la taille moyenne de l'index, en supposant que les 2/5� des docIds sont sur 1 octet, 2/5� sur 2 octets, et le reste en moyenne � 4 octets.
	\begin{correction}
		1 x 200 000 + 2 x 200 000 + 4 x 100 000= 1 Mo / vecteur.\\
		100 000 x 900k = 100 Go
	\end{correction}
\end{enumerate}
